\section{The State-Basis Discipline}
\label{sec:basis}
\label{sec:state-basis}

A corporate action rewrites the frame in which a unit's numbers mean anything. A 2-for-1
split turns 1{,}000 shares marked at 100 into 2{,}000 shares marked at 50: value is
unchanged, but every quantity and every price-like datum now speaks a different per-share
language. The ledger already records the entitlement moves --- ``a 2-for-1 split doubles
each entitled holding while the price reference halves'' (\S\ref{sec:contracts}) --- but
the second clause has no carrier: no stored fact says \emph{which frame} a consumed quote
inhabits, so a post-split balance can meet a pre-split quote and commit
$2{,}000\times100=\EUR{200{,}000}$, a phantom doubling both of whose factors are
individually correct. The defect is a missing referential-integrity edge between consumed
data and the state that gives them meaning; no internal reconciliation can catch it,
because nothing internal disagrees.

The repair is one coordinate and one seam. \emph{The coordinate is state; the datum
carries a reference; the seam is typed.} For each unit, the corporate-action basis is a
fold over logged events, materialised in \texttt{UnitStatus} behind the existing single
writer (\S\ref{sec:sec04-unitstatus}). Every consumed datum carries, in its attested
envelope, a stable reference into that fold. Quantity-like and price-like values meet
only inside a coherent snapshot whose basis assignment is fixed by the ledger and sealed
by a type variable; mixed-basis pairing is a type error at authorship and a typed refusal
at the erased store boundary. Informally: a datum is not a scalar but a pair (value,
state-version); pricing requires all inputs in the same state basis as the positions;
adjustment operators map data between state bases.

Observation time cannot serve as the coordinate. An external quote may lag the ledger's
own ex transition (\S\ref{sec:appD}); ex dates differ per venue; and a retro-effective
corporate action changes the basis of a past timestamp after the fact. A timestamp
therefore does not determine the frame, and a runtime join of two timelines would
re-derive at every consumption site what one state coordinate states once. Nor is a
fourth state home needed: observations already enter only as logged events
(\S\ref{sec:states}), and their store is a materialised projection of those events under
the existing cache discipline --- derivable, discardable, never authoritative. The
three-home model of \S\ref{sec:states} is retained; one field is added to the second
home.

\subsection{The basis coordinate}
\label{sec:basis-coordinate}

A unit's basis changes only at declared events. Two sources declare them, both
log-shaped.

\begin{definition}[Boundary event]
\label{def:boundary-event}
A \emph{boundary event} on unit $u$ is a logged event carrying the unit, an effective
time $t_{\mathrm{eff}}$, a \emph{declaration} --- a finite partial map from datum kinds
to operator specifications (\S\ref{sec:basis-operators}) --- and provenance. Boundary
events are content-addressed; \emph{bid} denotes the address. The two sources are
corporate-action lifecycle transactions on held units (\S\ref{sec:contracts}) and
external basis notices with no lifecycle leg --- index divisor maintenance, composition
true-ups --- entering as logged observation events with a canonical handler.
\end{definition}

\begin{definition}[Basis point, chain, effective order, epoch]
\label{def:basis-point}
A \emph{basis point} of $u$ is either the registration point $\bot_u$ or a boundary id
of $u$. The \emph{basis chain} of $u$, relative to a replay view, is the sequence of
$u$'s boundary events known to that view in \emph{effective order}: the lexicographic
total order on
\[
(t_{\mathrm{eff}},\ \mathit{prec},\ \mathit{bid}),
\]
where $\mathit{prec}$ is the intra-day precedence declared in the attested notice
(\S\ref{sec:basis-failure}, W4) and $\mathit{bid}$ is a deterministic final tie-break.
The \emph{epoch ordinal} of a basis point is its position in the chain --- a derived,
per-view quantity, never stored.
\end{definition}

Stamps and state store basis-point ids, never ordinals: a retro-effective boundary
inserts mid-chain and renumbers every later ordinal, while content addresses are stable.
The order is total on every representable log because the one case where the tie-break
is load-bearing --- two non-commuting boundaries on one unit with equal
$t_{\mathrm{eff}}$ --- is forced to carry declared precedence by the same-$t_{\mathrm{eff}}$
weld (W4): absent a declared mutual precedence, the colliding notice refuses into
pending-transition, fail-closed. Booking order never enters the composite; the tie-break
is data in the notice, and $\mathit{bid}$ is reached only where declared precedences tie.

\begin{definition}[The basis field]
\label{def:usbasis}
\texttt{UnitStatus} gains one field, \texttt{usBasis :: BasisPoint}, defined from
registration onward --- \texttt{defaultStatus} carries $\bot_u$ (\hyperref[cond:C5]{C5})
--- and written only through the closed \texttt{StatusWrite} set by one new constructor.
\end{definition}

\begin{lstlisting}[language=Haskell]
data BasisPoint = Origin UnitId | Boundary BoundaryId   -- content address, never an ordinal
  deriving (Eq, Ord, Show)

data StatusWrite                      -- the closed status-writer set, extended by ONE case
  = SetLifecycle    Lifecycle
  | SetLastSettle   (Price, BasisPoint)  -- was Maybe Qty: dimensionally wrong, basis-free
  | SetSupersededBy UnitId
  | SetBasis        BasisPoint        -- ABSOLUTE, not an increment: last-write-wins,
  deriving (Eq, Show)                 -- idempotent (P6), exactly as its siblings
\end{lstlisting}

\noindent
Re-applying \texttt{SetBasis b} to a status already at \texttt{b} is the identity, so P6
is untouched: the write is absolute, and the chain is a projection of the log --- no
counter is stored. The stated invariant is that \texttt{usBasis} equals the
effective-order tip of the chain as known to the view. This invariant is not
self-maintaining: replay folds the log in booking order (\S\ref{sec:lifecycle}), the
chain is projected in effective order, and the two folds do not commute on
retro-effective insertions --- under naked last-write-wins, a retro-effective boundary
booked late would regress the coordinate to a mid-chain point. The invariant is
therefore welded at the door, not assumed.

\begin{principle}[Tip weld]
\label{prin:tip-weld}
\texttt{applyTx} admits a transaction carrying \texttt{SetBasis}\,$b$ on unit $u$ only if
$b$ is the last element, in the effective order of
Definition~\ref{def:basis-point}, of $u$'s chain after insertion of the boundary the
transaction itself carries. The condition is checkable at admission from the committed
log and the transaction alone.
\end{principle}

For a chain-tip boundary --- the common case --- the mandated write is the boundary's own
id: behaviour unchanged. For a retro-effective boundary, the post-insertion tip is the
already-committed tip, so the mandated write re-asserts the standing tip --- idempotent
under P6 --- and the coordinate never moves backward: \emph{basis regressed by a late
notice} is an unrepresentable ledger state. The inserted boundary changes mid-chain
arrows, but arrows are per-view projections of the log, never stored
(Definition~\ref{def:basis-category}), so every subsequent snapshot recomputes them;
observations stamped in the window repair through re-stamp events
(\S\ref{sec:basis-time-travel}), moves through the existing compensating path. The
booking-order status fold and the effective-order chain projection consequently agree on
\texttt{usBasis} at every log prefix --- an equality witnessed by the retro-insertion
permutation oracle (Appendix~\ref{app:properties}), not merely asserted.

\begin{principle}[Invariance weld]
\label{prin:invariance-weld}
A basis-changing corporate action is one atomic transaction
(\hyperref[cond:C3]{C3}) carrying its entitlement moves, its \texttt{SetBasis}, and its
declaration; \texttt{applyTx} admits it only if the moves and the declaration jointly
satisfy the invariance identity of Theorem~\ref{thm:basis-value-invariance}: quantity
legs equal to $q\mapsto q\cdot f$ up to explicit cash-in-lieu legs, and booked cash legs
equal to $q\cdot c$.
\end{principle}

The check needs no market data and is the same admission idiom as \texttt{move}'s
positivity gate. Consequence: ``positions doubled but basis not advanced,'' ``basis
advanced but declaration missing,'' ``declaration inconsistent with moves,'' and --- by
the tip weld --- ``basis regressed by a late notice'' are all unrepresentable ledger
states.

\subsection{The adjustment operators and their composition law}
\label{sec:basis-operators}

An operator is declared data applied by one generic evaluator; no operator is inferred
from an event's name or from any model.

\begin{definition}[Operator specification]
\label{def:opspec}
A declaration maps each affected datum kind to one of the following. Parameters come
from the declared terms of the notice, never from the event taxonomy: a 5\% stock
dividend declaring $f=21/20$ is a scale regardless of its name.
\end{definition}

\begin{lstlisting}[language=Haskell]
data OpSpec
  = Scale Rational       -- q /= 0;  invertible
  | Shift Rational       -- invertible
  | Subst UnitId Rational -- r /= 0; cross-unit succession, invertible
  | Recompose Sigma      -- invertible iff sigma is declared bijective
  | AId                  -- identity: DECLARED, never assumed
  | Pending BoundaryRef  -- boundary known, parameter unpublished: no arrow yet
  | Terminal             -- series ends: no arrow in either direction
  deriving (Eq, Show)
-- Side conditions (non-zero, bijectivity) are enforced by unexported smart
-- constructors at the single admission door (applyTx), the `move` idiom: the
-- invalid spec is never stored, so the evaluator is total over stored values.
\end{lstlisting}

\noindent
The map is partial, and absence is load-bearing: a kind outside the declaration's domain
has no arrow and cannot cross that boundary --- fail-closed. \texttt{AId} must be
declared; it is never a default. All data-plane arithmetic is exact rational; the single
crossing to integer minor units is \texttt{toPrice}, whose rounding rule is fixed once in
\emph{ProductTerms} (round-half-even, one site), so composition of exact maps is
path-independent and intermediate caching is sound.

\begin{definition}[The basis category]
\label{def:basis-category}
Objects: the basis points of all units. Generating arrows, per datum kind $k$: for each
boundary event $g$ with $k$ in the domain of its declaration, one arrow from the point
before $g$ to $\mathit{bid}(g)$, with action the interpretation of the declared
operator; and for each declared succession (\texttt{Subst}, welded to
\texttt{usSupersededBy}, \hyperref[cond:C8]{C8}), one arrow from the predecessor's final
point to $\bot_{u'}$. Inverse arrows exist exactly where the generator is invertible;
\texttt{Pending} and \texttt{Terminal} generate no action. Within one unit the chain is
totally ordered and succession edges are unique, so between any two basis points there
is at most one arrow per kind.
\end{definition}

\begin{proposition}[Composition law]
\label{prop:composition-law}
For basis points $b\le b'$ on the (possibly succession-extended) chain, with intervening
boundaries $g_1,\dots,g_n$ in the effective order of Definition~\ref{def:basis-point},
\[
A_k(b\to b') \;=\; \mathrm{interp}(g_n)\circ\cdots\circ\mathrm{interp}(g_1),
\qquad\text{defined iff $k$ is in every $g_i$'s declaration;}
\]
\[
A_k(b'\to b) \;=\; A_k(b\to b')^{-1}
\quad\text{defined iff every generator is invertible;}
\]
\[
A_k(b\to b)=\mathrm{id},
\qquad
A_k(b\to b'') \;=\; A_k(b'\to b'')\circ A_k(b\to b').
\]
\end{proposition}

Operators do not commute; order-sensitivity is dissolved, not configured. Effective
order is data in the log --- $t_{\mathrm{eff}}$ and $\mathit{prec}$ from the attested
notice, $\mathit{bid}$ from content --- and no other composite is expressible, because
arrow values are authored solely by the ledger's fold: their constructors are unexported
(\S\ref{sec:basis-types}). Knowledge order lives on the bitemporal axis
(\S\ref{sec:basis-time-travel}), never in the composite.

The one illustration: spot 102 stamped at $b$; effective order dividend
\texttt{Shift 2} ($b\to b_1$), then split \texttt{Scale} $\tfrac12$ ($b_1\to b_2$).
Then $A(b\to b_2)(102)=(102-2)/2=50$. The reversed composite, $102/2-2=49$, is not an
arrow of the category: no pair of basis points has it as its unique arrow, and no
exported constructor can forge it. Had the dividend and split shared one
$t_{\mathrm{eff}}$, the composite would be fixed by their declared mutual precedence ---
and the pair would be inadmissible without it: the 49-versus-50 ambiguity is refused at
the door, never resolved by arrival order.

\begin{theorem}[Lifecycle value invariance across boundaries]
\label{thm:basis-value-invariance}
Let a boundary declare price action $p\mapsto(p-c)/f$ with quantity legs
$q\mapsto q\cdot f$ (up to explicit cash-in-lieu legs) and cash legs $q\cdot c$, as
Principle~\ref{prin:invariance-weld} requires. Then
\[
(q\cdot f)\cdot\frac{p-c}{f} \;=\; q\cdot p \;-\; q\cdot c :
\]
marked value is continuous through every boundary up to exactly the cash the same atomic
transaction books.
\end{theorem}

\begin{proof}
The identity is algebraic; the weld makes its hypotheses admission conditions, so every
committed boundary satisfies them. Fractional entitlements land as explicit cash-in-lieu
moves in the same transaction, so the identity holds to the minor unit; residues are
conserved flows, never rounding dust.
\end{proof}

\noindent
This proves, for basis boundaries, the lifecycle value-invariance property stated in
restricted form in \S\ref{sec:intro}: the price jump at the boundary exactly offsets the
booked cash.

\paragraph{Pending crystallisation.} A \texttt{Pending} boundary advances the basis
point at effectiveness --- the moves and \texttt{SetBasis} commit atomically
(\hyperref[cond:C3]{C3}) --- and its parameter arrives later as a
parameter-publication event referencing the boundary id, a versioned resolution under
the \emph{ProductTerms} append discipline (\hyperref[cond:C6]{C6}). The chain projection
resolves each boundary to its latest parameter version; no list element mutates, and the
epoch is unaffected because publication is not a boundary. Until publication, no arrow
crosses that boundary for the affected kinds: quarantine is forced
(\S\ref{sec:basis-failure}), which is correct, because any interim mark is a guess.

\begin{principle}[Re-derivation canon]
\label{prin:rederivation}
Primitive data are transported by declared arrows; derived data are re-derived from
transported inputs.
\end{principle}

The framework never assumes that a derivation commutes with adjustment: equivariance is
a property of the black box's interior, and the model-free directive forbids depending
on it. A stored derived datum --- a vendor divisor, an index level --- carries as its
stamp the joint basis of its inputs, a finite map $\mathrm{UnitId}\to\mathrm{BasisPoint}$
(a plain observation is the singleton case). It is admissible at a target assignment
$\beta$ iff its stamp equals $\beta$ pointwise on its whole domain
(Invariant~\ref{inv:basis}(ii)); the framework refuses to transport it, fail-closed. A
derivation performed inside a snapshot scope consumes transported primitive inputs and
is automatically at $\beta$ --- the phantom index propagating through ordinary
application.

\subsection{The stamp and the ingest door}
\label{sec:basis-ingest}

A raw vendor number becomes a ledger observation only through one total ingest function
--- the parse-boundary idiom of \texttt{move}.

\begin{lstlisting}[language=Haskell]
ingest :: Ledger -> SourceId -> Timestamp -> RawDatum
       -> Either IngestError StampedObs
data StampedObs   -- ABSTRACT: constructor unexported; ingest is the only door.
  -- Carries: value (exact), t_obs, source, and stamp :: Map UnitId BasisPoint,
  -- inside the attestation envelope (provider key, signature, content address).
\end{lstlisting}

\noindent
The stamp is assigned per (unit, source): a source lagging the ex transition
(\S\ref{sec:appD}) is a source still publishing at the old basis point while the unit
has advanced, recorded as a per-source basis offset --- either transcribed from the
source's declared convention under signature, or asserted by the gateway against the
committed corporate-action log. It is never defaulted from the ledger's prevailing
state. \texttt{ingest} additionally refuses a stamp naming an unregistered unit, so the
scope-closure condition of Invariant~\ref{inv:basis} is discharged at the door, never at
consumption. This subsumes the caller-asserted quote-ex flag and the hard-coded
one-step price adjustment of \S\ref{sec:appD}, generalising both to arbitrary boundary
chains; sourcing remains the market-data satellite's concern, while the stamp and the
coherence check are ledger-owned.

The assignment is a named trust assumption:

\begin{description}[style=nextline]
\item[TA-BASIS.] \emph{Owner:} market-data operations. \emph{Content:} the per-(unit,
source) basis stamp assigned at ingest is true of the source's dissemination.
\emph{Violation consequence:} mis-based valuation. \emph{Detection:} W3 partition
quarantine (\S\ref{sec:basis-failure}) --- a large innovation coinciding with no
committed boundary quarantines the series --- and daily reconciliation of
vendor-published adjustment factors against applied declarations.
\end{description}

\subsection{The invariant}
\label{sec:basis-invariant}

\begin{invariant}[Single-basis consumption]
\label{inv:basis}
Let $V$ be a valuation invocation over position scope $S$ at time $t$, and let
$\beta_t:\mathrm{UnitId}\to\mathrm{BasisPoint}$ be the total basis assignment given by
replayed \texttt{usBasis} over all registered units (total by
Definition~\ref{def:usbasis}: \texttt{defaultStatus} carries $\bot_u$). Then:
(i)~every balance consumed by $V$ is read from $S$ at $\beta_t|_S$;
(ii)~every datum $d$ consumed by $V$ satisfies
$\mathrm{stamp}(d)(u)=\beta_t(u)$ for every $u\in\mathrm{dom}(\mathrm{stamp}(d))$ ---
the check ranges over the stamp's whole domain, whether or not $u\in S$;
(iii)~every value formed from several data is formed from data satisfying (ii) under one
and the same $\beta_t$.
\end{invariant}

\noindent
\emph{Well-definedness.} \texttt{withSnapshot} (\S\ref{sec:basis-types}) fixes $\beta$
as $\beta_t$ restricted to the \emph{stamp-closure} of $S$ ---
$S\cup\bigcup\{\mathrm{dom}(\mathrm{stamp}(d)) : d \text{ consumed}\}$ --- which is
defined because $\beta_t$ is total over registered units and \texttt{ingest} admits no
stamp naming an unregistered unit. $S$ restricts which balances are read; it never
restricts the domain on which admissibility is checked.

Clause (iii) is what makes the invariant meaningful for a black box: the framework
legislates which data it is legal to feed a pricing function, never what the function
does. State-sufficiency (\S\ref{sec:valuation}) survives with its third, previously
silent dependency --- basis agreement --- now carried by the snapshot's type variable;
the three illegal states of \S\ref{sec:valuation} gain a fourth: a held unit priced in
the wrong basis is representable as an error and unrepresentable as a \texttt{Cash}.

The condition form, alongside C1--C12 (\S\ref{sec:sec04-conditions}):

\begin{description}[style=nextline]
\item[C13 (basis edge).]\hypertarget{cond:C13}{}\label{cond:C13} Every consumed datum
names its joint basis; consumption requires pointwise agreement, over the stamp's whole
domain, with the prevailing total assignment $\beta_t$, or a ledger-authored arrow. C13
is the data-plane sibling of P3: referential integrity for meaning, as P3 is for
existence.
\end{description}

\subsection{Type-level enforcement}
\label{sec:basis-types}

The guarantee is stated honestly, in the two-tier discipline of
\hyperref[cond:C11]{C11}: typed at authorship, erased at rest, re-established at
consumption. Epochs are runtime replay outputs, so the honest static guarantee is
\emph{sameness}, not value --- a type-level ordinal in the normative layer would claim
static knowledge the system cannot have.

\paragraph{Tier 1 --- the didactic law.} One page of exposition of the invariant, not
the system mechanism. All constructors are unexported; the sole author is the ledger's
fold, so an \texttt{Adj} value \emph{is} a proof that the log carries exactly these
boundaries, in effective order.

\begin{lstlisting}[language=Haskell]
{-# LANGUAGE DataKinds, GADTs, KindSignatures, RoleAnnotations #-}
data N = Z | S N

newtype PriceAt (e :: N) = PriceAt Rational   -- constructor NOT exported
newtype BalAt   (e :: N) = BalAt   Integer    -- constructor NOT exported
type role PriceAt nominal   -- GHC would infer PHANTOM, and Data.Coerce.coerce
type role BalAt   nominal   -- (Safe Haskell, no unsafe feature) would then forge
                            -- the crossing at ANY pair of indices. Nominal makes
                            -- the coercion derivable only at equal index.

data Adj (m :: N) (n :: N) where              -- constructors NOT exported:
  AIdA  :: Adj n n                            -- the ledger fold is the only author
  AStep :: Adj m n -> OpSpec -> Adj m ('S n)  -- one boundary, in effective order

adjust :: Adj m n -> PriceAt m -> PriceAt n   -- the ONLY epoch-crossing function

markValue :: BalAt e -> PriceAt e -> Cash     -- Qty and Price meet ONLY at equal index
markValue (BalAt q) (PriceAt p) = ...

-- _bad :: BalAt ('S e) -> PriceAt e -> Cash
-- _bad = markValue
--   TYPE ERROR. The epoch-(n+1) position priced with the epoch-n quote is not a
--   wrong number; it is not a program. And AStep extends only to 'S n, so the
--   mis-ordered composite of the 49-vs-50 illustration has no inhabitant.
\end{lstlisting}

\paragraph{Tier 2 --- the normative seam.} \texttt{withSnapshot} is the one constructor
of coherent read scopes, mirroring \texttt{applyTx} as the one write door. The ledger is
sealed inside the scope --- no \texttt{Ledger} parameter escapes downstream; the skolem
\texttt{b} denotes a whole basis assignment, not one unit's epoch, and every accessor is
unit-keyed; the per-unit price accessor is fallible, so quarantine is representable.

\begin{lstlisting}[language=Haskell]
data Snapshot b   -- ABSTRACT. Constructed only by withSnapshot, which fixes the
                  -- target assignment beta = beta_t on the STAMP-CLOSURE of the
                  -- scope (Invariant B: beta_t is total over registered units;
                  -- Scope restricts which balances are read, never beta's domain),
                  -- reads balances from the SAME sealed ledger, and transports
                  -- every admissible stamped observation to beta along declared
                  -- arrows -- refusing, per unit and kind, where none exists.
type role Snapshot nominal

newtype BalAt'   b = ...        -- constructors NOT exported
newtype PriceAt' b = ...
newtype At b a     = ...
type role BalAt'   nominal
type role PriceAt' nominal
type role At       nominal nominal

withSnapshot :: Ledger -> Timestamp -> Scope
             -> (forall b. Snapshot b -> r) -> r     -- quantifier where the proof lives

snapBal   :: Snapshot b -> WalletId -> UnitId -> BalAt' b
snapPrice :: Snapshot b -> UnitId -> Either BasisError (PriceAt' b)  -- per-unit refusal

markV :: BalAt' b -> PriceAt' b -> Cash                  -- the seam; Cash is basis-free
zipAt :: (x -> y -> z) -> At b x -> At b y -> At b z     -- in-scope derivation

data Valuation = Priced Cash
               | Unpriced UnitId BasisError              -- quarantined, enumerated
               | PricedStale UnitId BasisPoint BasisPoint Cash  -- logged election

value :: Snapshot b -> [WalletId] -> (Cash, [Valuation]) -- book total + exclusions

-- The single cross-scope door, ledger-authored, factoring ONLY through declared
-- arrows:
carry :: Snapshot b0 -> Snapshot b1 -> Kind a -> At b0 a
      -> Either BasisError (At b1 a)
\end{lstlisting}

\paragraph{Unreachability, for an arbitrary black-box pricing function.} Let
$F::\forall b.\ \texttt{Snapshot } b\to r$ be any pricing function. By parametricity in
\texttt{b} --- within the fragment stated below --- every \texttt{b}-indexed value $F$
manipulates derives from its one \texttt{Snapshot} argument: the carriers' constructors
are unexported, their only introduction forms are \texttt{snapBal}, \texttt{snapPrice},
\texttt{carry}, and index-preserving combinators, and each yields values at $\beta$ by
construction of the door. Two snapshots' skolems never unify, and \texttt{carry} --- the
only cross-skolem function --- factors through the unique declared arrow or refuses.
Hence every (quantity, price-like) pair any well-typed $F$ can form lies in a single
basis assignment. The runtime complement, at the erased store boundary, is the
witness-equality check inside \texttt{withSnapshot}, which returns a typed
\texttt{BasisError}, never a silent pick. This is the C11 guarantee restated at the read
seam: compile-time at authorship, typed refusal at the boundary --- claimed as exactly
that and no more.

\paragraph{Language-fragment proviso.} GHC infers role \emph{phantom} for a type
parameter occurring in no constructor field; \texttt{Coercible} then holds at all index
pairs, and \texttt{Data.Coerce.coerce} --- Safe Haskell, no unsafe feature --- would
forge the epoch and skolem crossings past every unexported constructor. Every
basis-indexed carrier therefore declares \texttt{type role \dots\ nominal} on its basis
index; the annotations are normative, not stylistic. The enumeration above is claimed
for exactly this fragment: Safe Haskell; nominal roles on every basis-indexed parameter;
constructors, field selectors, and eliminators unexported; no \texttt{Generic},
\texttt{Data}, or \texttt{Typeable}-derived path exposing representation; no
\texttt{GeneralizedNewtypeDeriving} on the carriers; no Template Haskell in consuming
modules. Within the fragment the free-theorem reading of $\forall b$ is sound and
\texttt{coerce} contributes no introduction form.

\subsection{Time travel and late corporate actions}
\label{sec:basis-time-travel}

\texttt{clone\_at($t$)} already rebuilds every \texttt{UnitStatus} field by the
catamorphism (\S\ref{sec:lifecycle}); \texttt{usBasis} is such a field, so the basis
prevailing at $t$ is reconstructed with zero new machinery (P8, extended). Historical
valuation is \texttt{withSnapshot} over the clone: the target assignment is $\beta_t$;
raw observations are immutable and stamped, so they never migrate; a datum stamped after
$\beta_t$ is transported backward iff every intervening generator is invertible
(Definition~\ref{def:opspec}), else refused into the \S\ref{sec:basis-failure} workflow.

Two version axes are orthogonal, and both compose. The \emph{basis axis} is
ledger-authored: effective order, per Proposition~\ref{prop:composition-law}. The
\emph{knowledge axis} is bitemporal: ``as known at $t$'' replays with the chain as known
at $t$; ``restated'' replays with corrections. A late or retro-effective boundary
(booking time $t_{\mathrm{notify}}$, economic time $t_{\mathrm{eff}}$) changes no stored
stamp silently: stamps are stable ids, and the ledger appends \emph{re-stamp derivation
events} for observations captured in $[t_{\mathrm{eff}},t_{\mathrm{notify}})$ --- value
unchanged, coordinate corrected, lineage to the late transaction. Its \texttt{SetBasis}
is tip-welded (Principle~\ref{prin:tip-weld}): it re-asserts the standing effective tip,
so the live coordinate never regresses; the corrected replay's chain projection places
the boundary mid-chain and recomputes the spanning arrows. The as-known replay
reproduces the honest mistake; the corrected replay applies the re-stamps; moves
committed in the window repair through the existing compensating-transaction path
(\S\ref{sec:sec11-corrections}), now with a citable coordinate lineage. A pinned
snapshot plus a pinned chain version is reproducible \emph{and} right; pinning alone was
reproducibly wrong.

\subsection{Failure semantics as workflow}
\label{sec:basis-failure}

Every refusal is typed (\texttt{BasisError}); consequence class decides the regime;
nothing is silent anywhere.

\renewcommand{\arraystretch}{1.25}
\begin{center}
\begin{longtable}{@{}p{2.2cm} p{3.6cm} p{4.7cm} p{2.9cm}@{}}
\toprule
\textbf{Regime} & \textbf{Trigger} & \textbf{Consequence} & \textbf{Mandatory for} \\
\midrule
\endhead
W1 Block & Move-emitting consumption (settlement projection, NAV strike, fee
crystallisation) meets any \texttt{BasisError}. & The lifecycle event defers as a
recorded workflow event with a named unblock condition --- a fresh observation at
$\beta$, or the crystallising publication event --- queued under obligation liveness
(P21, \S\ref{sec:obligation-liveness}). Real money never moves on an unwitnessed basis.
& The stale-data gate (\S\ref{sec:sec11-faults}) and the pre-invocation check
(\S\ref{sec:temporal}): basis equality first, timestamp threshold second. \\
\addlinespace
W2 Per-unit quarantine & \texttt{snapPrice} returns \texttt{Left} for a unit: no
admissible datum, a \texttt{Pending} gap, an undeclared kind. & The book values on time
as $(\texttt{Cash},[\texttt{Unpriced}\ \dots])$ --- a total minus a named, enumerated
exclusion list. Blocking the whole book is rejected: one late notice must not halt
thousands of units; shippability decides. & Daily marks, risk, management PnL. \\
\addlinespace
W2$'$ Flagged-stale carry & An explicit, logged election event with its canonical
writer row in the C11 table. & The unit is valued \emph{coherent-at-}$\beta^-$ --- the
latest assignment $\beta^-\le\beta_t$ at which its data are pointwise complete --- never
a mixed vector. The result is \texttt{PricedStale}, a distinct constructor carrying its
gap; downstream must pattern-match. \emph{Coherent-at-}$\beta$ means all inputs share
assignment $\beta$; \emph{current} means $\beta=\beta_t$; W2$'$ is coherent, not
current, and says so. & Management PnL where quarantining a large unit is worse than a
flagged carry. \\
\addlinespace
W3 Partition quarantine & Aggregation partitions sources by asserted basis point;
quorum and divergence are computed within a partition only. & $\{100\ @\ b_3\}$ and
$\{50\ @\ b_4\}$ are two singleton partitions --- no false divergence, no median 75 that
exists in no basis. A large innovation with no committed boundary is the signature of an
undeclared corporate action or a bad feed: the series quarantines. This is TA-BASIS's
detection signal; the window in which the world has split but the log has not is
detected, not prevented. & All ingestion around corporate-action windows. \\
\addlinespace
W4 Notice attestation & A boundary declaration rewrites both sides of the seam and,
once its moves are emitted, cannot be repaired by the price-correction path. & Admission
requires issuer or exchange signature, or content-hash agreement across
$N_{\mathrm{CA}}\ge2$ independent sources (owner: data governance). Same-$t_{\mathrm{eff}}$
weld: a boundary whose $t_{\mathrm{eff}}$ equals that of another boundary on the same
unit --- committed or co-presented --- is admitted only if the attested notices declare
mutual precedence ($\mathit{prec}$, Definition~\ref{def:basis-point}); absent
declaration, the colliding notice refuses into pending-transition, fail-closed. Between
first notice and confirmation the unit is \emph{pending-transition}: W1 blocks, W2
quarantines, W3 partitions. & Every boundary event. \\
\bottomrule
\end{longtable}
\end{center}

The pending spin-off is the normal case, not the edge: \texttt{Pending} advances the
basis point at effectiveness, so no number prices the unit at the new basis until the
parameter publishes --- quarantine is forced, and a flagged carry at the old basis is
the only honest alternative.

\subsection{The 2-for-1 split, end to end}
\label{sec:basis-worked}

A wallet holds 1{,}000 shares of $u$; $\texttt{usBasis}=b_3$.

\begin{enumerate}[nosep]
\item \emph{Snap.} Vendor EX1 publishes 100.00; \texttt{ingest} stamps
$\{u\mapsto b_3\}$ per the $(u,\mathrm{EX1})$ convention. Immutable, logged.
\item \emph{The corporate action.} One atomic transaction $\tau$: moves $+1{,}000$ from
the corporate-action virtual wallet ($1{,}000\to2{,}000$); \texttt{SetBasis}\,$b_4$;
declaration $\{\mathrm{KSpot}\mapsto\texttt{Scale}\ \tfrac12,\dots\}$. Admission welds:
quantity legs $2{,}000=1{,}000\times2$; cash legs $0=1{,}000\times0$
(Theorem~\ref{thm:basis-value-invariance}); $b_4$ is the post-insertion effective tip
(Principle~\ref{prin:tip-weld}). No reachable state has the doubling without the
boundary.
\item \emph{Snapshot at $t_{\mathrm{price}}$.} $\beta(u)=b_4$; balance 2{,}000 at $b_4$
by the fold; the stored spot is stamped $b_3$; the door transports along the unique
arrow: $A(b_3\to b_4)(100)=50$.
\item \emph{Mark.} $\texttt{markV}\ 2{,}000\ 50=\EUR{100{,}000}$. Phantom: none.
\item \emph{The bug, attempted.} Feeding the raw 100 requires a
$\texttt{PriceAt'}\ b$ value, whose only doors are \texttt{snapPrice} and \texttt{carry}
--- both refuse or transport --- and whose nominal role bars the \texttt{coerce} back
door. In Tier 1, pairing the epoch-$(n{+}1)$ balance with the epoch-$n$ quote is a type
error. The \EUR{200{,}000} is a program that does not exist on the compile side and a
\texttt{Left BasisMismatch} on the runtime side; it never reaches a committed
\texttt{Cash}.
\item \emph{Attribution.} True post-split spot 55 (up 10\% like-for-like). Open
$\texttt{Snapshot}\ b_0$ (pre) and $\texttt{Snapshot}\ b_1$ (post); \texttt{carry}
transports the opening pair along the declared arrow, $(1{,}000,100)\mapsto(2{,}000,50)$,
value-preserving by Theorem~\ref{thm:basis-value-invariance} with $c=0$. Then
$\mathrm{PnL}_{\text{price}}=2{,}000\times(55-50)=+10{,}000$ and
$\mathrm{PnL}_{\text{flow}}=(2{,}000-2{,}000)\times55=0$. The $-45{,}000/+45{,}000$
phantom pair of the naive decomposition is gone: the decomposition, not just the sum, is
basis-invariant, and it is writable through the one declared door.
\item \emph{Late notice.} If $\tau$ is unlogged at $t_{\mathrm{price}}$,
\hyperref[cond:C3]{C3} keeps moves and boundary together, so the snapshot is coherent at
$b_3$: $1{,}000\times100=\EUR{100{,}000}$ --- right under the information held.
Meanwhile the exchange disseminates $\approx50$: W3 fires and quarantines the series; W1
blocks settlement-grade use. At $t_{\mathrm{notify}}$, $\tau$ commits with economic time
$t_{\mathrm{eff}}$; if later-effective boundaries committed in the interim, the tip weld
makes $\tau$'s \texttt{SetBasis} re-assert the standing tip while its boundary inserts
mid-chain in the projection. Re-stamp events move the window's observations to their
corrected points; the corrected replay is coherent end to end. At no point does
$2{,}000\times100$ or $1{,}000\times50$ reach a committed \texttt{Cash}.
\end{enumerate}

\paragraph{The dividend.} The same walk with \texttt{Shift 2} and $f=1$: transport gives
$102-2=100$; the 2-per-share cash leg books once, in the same transaction, with
$\text{booked cash}=q\cdot c$ by the weld; total \EUR{102{,}000}. The
\texttt{Cum}/\texttt{Ex} pricing of \S\ref{sec:appD} is recovered exactly as the
one-step, $f=1$ fibre: \texttt{Cum} is the old basis point, \texttt{Ex}~$d$ the new one
with generator \texttt{Shift}~$d$, and the quote-ex flag is the per-(unit, source)
stamp. The lagging source --- quote still cum after the ledger's ex transition --- is
stamped at the old point and transported: the case \S\ref{sec:appD} hard-codes,
generalised.

\paragraph{The index divisor.} Index $I$ over A (60) and B (84); level $(60+84)/D=100$
fixes $D=1.44$, snapped as a derived observation with joint stamp
$\{I\mapsto b_I,\ A\mapsto b_A,\ B\mapsto b_B\}$. B splits 2-for-1 ($b_B\to b_B'$);
fresh prices A 60, B 42. The corrupt number --- $(60+42)/1.44=70.8$, a phantom
$-29.2\%$ index move --- requires consuming the snapped $D$ at an assignment where
$\beta_t(B)=b_B'\ne b_B$: pointwise inadmissible, refused fail-closed --- and the check
is well-defined although the scope holds only $I$, because $\beta_t$ is total over
registered units and admissibility ranges over the stamp's full domain $\{I,A,B\}$
(Invariant~\ref{inv:basis}(ii)). The framework will not transport a derived datum
(Principle~\ref{prin:rederivation}). Legal paths: the index authority's maintenance
boundary on $I$ declares \texttt{Recompose} with $D'=1.02$, so $(60+42)/1.02=100$,
entering under W4 like any corporate action; a fresh divisor observation arrives stamped
at the new assignment; or the level is re-derived in-scope from transported
constituents. Until one of these, $\texttt{snapPrice}\ I$ is a \texttt{Left}:
\texttt{Unpriced}~$I$, enumerated, never a silent 70.8.

\subsection{CDM corporate-action mapping}
\label{sec:basis-cdm}

CDM \texttt{BusinessEvent}s already map to transactions through the forgetful map $F$
(\S\ref{sec:cdm}); ``CDM'' denotes Common Domain Model v6.0.0 throughout, as pinned in
Appendix~\ref{sec:appF}. The extension is one rule: a corporate-action event's
instructions supply both halves of $\tau$ --- the moves and the declaration --- keyed on
declared terms, never taxonomy. $F$ lands in (moves, \texttt{SetBasis} + declaration)
instead of moves alone; it remains total and deterministic, and the one economic fact it
now keeps is precisely the fact it was discarding.

\renewcommand{\arraystretch}{1.25}
\begin{center}
\begin{longtable}{@{}p{4.4cm} p{4.2cm} p{4.8cm}@{}}
\toprule
\textbf{CDM representation} & \textbf{Moves (existing)} & \textbf{Declaration (new)} \\
\midrule
\endhead
StockSplit / ReverseSplit / StockDividend (QuantityChange, ratio $r$) & entitlement
$\times r$ via the corporate-action virtual wallet (\S\ref{sec:contracts}) &
$\mathrm{KSpot}\mapsto\texttt{Scale}\,(1/r)$, per-share reference kinds likewise; weld
checks $q'=q\cdot r$ \\
\addlinespace
CashDividend / coupon detachment (Transfer) & cash issuer $\to$ holders &
$\mathrm{KSpot}\mapsto\texttt{Shift}\,d$ on cum-quoted kinds; weld checks cash
$=q\cdot d$; subsumes \S\ref{sec:appD} \\
\addlinespace
SpinOff & issuance moves of the child; child registered at $\bot$ & parent kinds
\texttt{Pending}\,\textit{ref} until allocation publishes; the publication event
crystallises the parameter \\
\addlinespace
Merger, stock-for-stock (TermsChange + succession) & C8 Breaking amendment + paired
issuance; \texttt{usSupersededBy} & \texttt{Subst}\,$u'\,r$ --- the cross-unit arrow;
predecessor data reach the successor's basis through it \\
\addlinespace
Merger for cash & terminal moves & \texttt{Terminal} --- no arrow in either direction;
fail-closed forbids any further crossing \\
\addlinespace
Index rebalance / divisor change & rebalance moves on the strategy unit & boundary on
the index unit via the basis-notice handler;
$\mathrm{KComposition}\mapsto\texttt{Recompose}\,\sigma$, under W4 \\
\addlinespace
Election / proration windows & per-holder entitlement differences are moves & boundary
\texttt{Pending} during the window; the unit-level operator crystallises at proration
--- Status stays read identically by every holder \\
\addlinespace
TermsChange (non-corporate-action amendment) & C6/C8 as today & no declaration; the
basis advances only on declared basis change \\
\bottomrule
\end{longtable}
\end{center}

\subsection*{Key Invariants and Consequences}

\begin{itemize}[nosep]
\item \textbf{One coordinate, one seam (C13).} The corporate-action basis is a
\texttt{UnitStatus} field folded from logged boundary events; every consumed datum
carries a joint basis stamp; consumption requires pointwise agreement with the total
assignment $\beta_t$ or a ledger-authored arrow. C13 is the data-plane sibling of P3.
(\S\ref{sec:basis-coordinate}, \S\ref{sec:basis-invariant}.)
\item \textbf{Value invariance is a theorem, not a property claim.} The invariance weld
makes the identity $(qf)\cdot\frac{p-c}{f}=qp-qc$ an admission condition, so marked
value is continuous through every boundary up to exactly the booked cash
(Theorem~\ref{thm:basis-value-invariance}); the lifecycle value-invariance property of
\S\ref{sec:intro} is proved on both factors for basis boundaries.
\item \textbf{No regression, no reordering.} \texttt{SetBasis} is absolute and
idempotent, so P6 is unchanged; the tip weld makes ``basis regressed by a late notice''
unrepresentable, and booking-order replay agrees with the effective-order chain
projection on \texttt{usBasis} at every log prefix --- witnessed by the retro-insertion
permutation oracle (P24 family, Appendix~\ref{app:properties}), not asserted.
\item \textbf{Fail-closed by expressibility.} Declarations are partial maps: an
undeclared kind has no arrow and cannot cross a boundary; \texttt{Pending} and
\texttt{Terminal} generate no action; identity is declared, never assumed.
(\S\ref{sec:basis-operators}.)
\item \textbf{The phantom is not a number but a non-program.} Within the stated language
fragment, every (quantity, price-like) pair a well-typed pricing function can form lies
in one basis assignment; at the erased boundary the same discipline is a typed refusal,
never a silent pick. Mixed-basis value is unrepresentable as a \texttt{Cash} and
representable only as an enumerated \texttt{Valuation} exclusion.
(\S\ref{sec:basis-types}, \S\ref{sec:basis-failure}.)
\item \textbf{Time travel extends for free.} \texttt{clone\_at($t$)} rebuilds
\texttt{usBasis} by the existing catamorphism (P8); the basis axis is effective order,
the knowledge axis is bitemporal, and late notices repair through re-stamp events with
citable lineage while the live coordinate never moves backward.
(\S\ref{sec:basis-time-travel}.)
\item \textbf{Protected elements untouched.} The atomic move, conservation, and log
immutability (P1, P2, P4) are unchanged in obligation; every addition is an appended
event class, a new \texttt{StatusWrite} constructor, a new admission condition at the
existing single door, a named projection, or a type refinement at the seam.
\end{itemize}

\noindent
The pairing is the invariant: type both of its factors in the same frame, and the
phantom PnL has nowhere to live.
